\subsection{Mixed first coordinates at three and an inert prime}

\begin{statement}{E.55}{Proved}{A free-product arithmetic presentation and every mixed first-coordinate fiber}
Let $\ell$ be a rational prime congruent to $2$ or $5$ modulo
$9$, and put $q=\ell^2$. In particular $\ell=2$ is allowed.
Retain the actual based arithmetic quotient $H$ of E.46.
For $a\in\mathbb Z_3\setminus\{0\}$ and $b\in\mathbb Z_3$,
fix the first Kummer class
\[
 [A]=[3a\kappa_3(\ell)+3b\kappa_3(3)],\qquad
 A_0=3a=3^s\epsilon,\quad s=v_3(3a),\quad
 \epsilon\in\mathbb Z_3^\times,\quad B_0=3b.
 \tag{E55.1}
\]
For a rational prime $p\notin\{\ell,3\}$, set $h_p=v_3(p-1)$.
Choose integer representatives $a_p,b_p$ of $a,b$ modulo
$3^{\max(h_p-1,0)}$, and define
\[
 \mathscr P_{a,b}^{\ell}=
 \left\{p\notin\{\ell,3\}:
 \left(\ell^{3a_p}3^{3b_p}\right)^{(p-1)/3^{h_p}}
                         \equiv1\pmod p\right\}.
 \tag{E55.2}
\]
This condition is independent of the representatives. Write
$\mathcal R(\mathscr P)$ for the restricted three-adic product
defined in (E49.2). After the explicit arithmetic choices below,
the complete global section fiber of $H$ over $[A]$ has additive
topological coordinates
\[
 \boxed{\quad
 \mathbb Z_3^2\times\mathcal R(\mathscr P_{a,b}^{\ell})
                    \times\mathbb Z_3^s\times\mathbb Z_3[[Z]].
 \quad}                                                     \tag{E55.3}
\]
Writing the coordinates as
$(\beta_\ell,\beta_3,(\beta_p),(z_d)_{d=1}^s,f)$,
the scaled toric map is exactly
\[
 \left(a\kappa_3(\ell)+b\kappa_3(3),\quad
 \beta_\ell\kappa_3(\ell)+\beta_3\kappa_3(3)
                +\sum_{p\in\mathscr P_{a,b}^{\ell}}
                                      \beta_p\kappa_3(p)\right).
 \tag{E55.4}
\]
Every fiber over an allowed second coordinate is therefore
$\mathbb Z_3^s\times\mathbb Z_3[[Z]]$, with the displayed
integral coordinate maps. The classes with $a=0$ are those
computed in E.47 and E.51. Thus these statements together
compute every first class supported on $\{\ell,3\}$.

The arithmetic input is an actual continuous pro-three
isomorphism. With $K=\mathbb Q(\zeta_3)$ and
$S=\{\text{primes of }K\text{ over }\ell,3\}$, there are
generators $t,\sigma,\rho$ for which
\[
 G_{K,S}^{(3)}\simeq
 \langle t,\sigma\mid\sigma t\sigma^{-1}=t^q\rangle_{\!3}
                  \mathop{\ast}_{3}\langle\rho\rangle,
 \qquad \langle\rho\rangle\simeq\mathbb Z_3.                 \tag{E55.5}
\]
Here $t$ is actual tame inertia, and $\sigma$ is actual
arithmetic Frobenius at the unique $K$-prime over $\ell$.
In the relative pro-three quotient over
$\operatorname{Gal}(K/\mathbb Q)=\langle c\rangle$ one has
\[
 \begin{array}{c|ccc}
       &t&\sigma&\rho\\\hline
 \chi_3&1&q&1\\
 \kappa_3(\ell)&1&0&0\\
 \kappa_3(3)&0&0&1\\
 c(\,\cdot\,)c&t^{-1}&\sigma&\rho^{-1}.
 \end{array}                                                \tag{E55.6}
\]
The projection from the full relative arithmetic group onto
this group has a continuous $c$-equivariant section retaining
the actual decomposition group at $\ell$.

For the part unramified outside $\{\ell,3\}$, all second
cocycles have unique representatives with $v(c)=0$ and
\[
 \begin{split}
 u=v(t)&=T^{A_0/2}\left(3\beta_\ell+
            \sum_{d=1}^s z_d\bigl(Q_d(T)-\varphi(3^d)\bigr)\right),\\
 Q_d(T)&=\sum_{\substack{-(3^d-1)/2\le j\le(3^d-1)/2\\3\nmid j}}
                     T^{\epsilon3^{s-d}j},\\
 v(\sigma)&=
 \frac{q\,u(T^q)-(1+T^{A_0}+\cdots+T^{(q-1)A_0})u(T)}
                                      {T^{qA_0}-1},\\
 v(\rho)&=T^{B_0/2}\bigl(3\beta_3+Vf(V)\bigr),
                    \qquad V=T+T^{-1}-2.
 \end{split}                                                \tag{E55.7}
\]
The quotient in this formula belongs to $\mathbb Z_3[[T-1]]$,
is anti-invariant under inversion, and has augmentation zero.
No division by a nonunit is left unspecified.
The proof constructs the extra global classes by actual
homomorphisms on the arithmetic inertia kernel and gives their
precise toric normalization.

The program \texttt{e\_two\_prime\_first.py} verifies the
coordinate equations and integral divisions, recovers arbitrary
tested representatives from cyclic group rings, and checks the
residue criterion by four independent finite descriptions.

\emph{Manuscript: \S51.16, folios733v--736v, scans172--178,
(371)--(379), Part E(c),(d) and the fourth exactness follow-up.}
\end{statement}
\begin{proof}
\emph{The actual arithmetic generators.}
All groups in this paragraph are relative pro-three Galois
groups over $K$, together with the indicated order-two action.
The local group at $\ell$ has tame inertia $\mathbb Z_3(1)$
and unramified quotient $\mathbb Z_3\times C_2$ over
$\mathbb Q_\ell$. Its injection into the full relative global
group follows from the inertia injection in E.39 and the
injective unramified cyclotomic character. Write its Frobenius
as $F$. The order-two component $d$ of its procyclic unramified
subgroup acts on inertia by inversion. The element
$\sigma=F^2$ commutes with $d$ and acts on inertia by $q$.
Choose the Frobenius complement to fix compatible roots of the
uniformizer $\ell$. Cubing is an automorphism of
$\mathbb Z_\ell^\times$: its residue group has order
$\ell-1$ prime to three, and its principal units are pro-$\ell$.
Consequently the compatible roots of $3$ are in $\mathbb Q_\ell$.
With these root choices the two Kummer coordinates vanish on
$d$ and $\sigma$.

The finite Schur--Zassenhaus conjugacy theorem, followed by
compact inverse limits, conjugates the image of a complex
conjugation to $d$. We use this conjugate and denote it by $c$.
Choose the compatible roots of positive rational primes fixed
by this $c$. The roots of $\ell$ and $3$ just chosen are such
roots. This is a permitted change of representatives of the
same rational Kummer classes. More explicitly, replacing roots
by $\zeta_{3^n}^{e_i}$ times the old roots changes their Kummer
cocycles by $(\chi_3-1)e_i$. For (E55.1) the change is
$(\chi_3-1)3(ae_\ell+be_3)$; conjugation by the element
$(-3(ae_\ell+be_3),0)\in H$ effects exactly this change of first
cocycle, using the multiplication law of E.46. Thus the based
section fibers are transported by an explicit geometric
conjugation. The coefficient involution after normalization is
still $C(f)=-f(T^{-1})$.

Put $P=G_{K,\{3\}}^{(3)}$. E.34 and E.37 prove that $P$ is
free pro-three of rank two. The characters given by the
cyclotomic character modulo nine and by $\kappa_3(3)$ modulo
three are independent on its Frattini quotient. Since
$v_3(q-1)=1$, the image of $\sigma$ is primitive in that
quotient, has nonzero first coordinate and zero second
coordinate. Choose an element $\gamma$ of the full group whose
image has second coordinate a unit. Multiply it by
$t^{-\kappa_3(\ell)(\gamma)}$ to obtain $\gamma_0$ with
$\kappa_3(\ell)(\gamma_0)=0$. The exact element
\[
 r=\gamma_0c\gamma_0^{-1}c
 \quad\text{satisfies}\quad
 crc=r^{-1},\quad\chi_3(r)=1,\quad
 \kappa_3(\ell)(r)=0,\quad
 \kappa_3(3)(r)=2\kappa_3(3)(\gamma_0).
 \tag{E55.8}
\]
Raise $r$ to the inverse of the last, unit, coordinate to
obtain $\rho$. Its image and $\sigma$ are a free basis of $P$:
their Frattini images form a basis, and the resulting
surjective endomorphism of a free pro-three group of finite
rank is an automorphism. This proves all the Kummer and
involution identities in (E55.6) on actual arithmetic elements.

\emph{The free-product comparison and its section.}
Apply the opened arithmetic Riemann-existence theorem
\cite[Theorem2.4]{E_WingbergRiemann} to $\{3\}\subset S$.
Its hypotheses are exactly those checked in E.39: the base
has the $K(\pi,1)$ property, and the cyclotomic tower gives
nontrivial unramified completion at $\ell$. Its kernel is the
continuous free pro-three product of the inertia groups over
the coset space $P/\langle\sigma\rangle$.
The actual decomposition group
$D=\langle t,\sigma\mid\sigma t\sigma^{-1}=t^q\rangle_3$
and $\rho$ therefore induce a continuous surjection
$D\ast_3\langle\rho\rangle\to G_{K,S}^{(3)}$ over $P$.

Its map on kernels is an isomorphism. One can see the kernel
of the source directly in its pro-three coset tree: the
vertices are the cosets of $D$ and $\langle\rho\rangle$,
and the edges are the group itself. Quotienting by that kernel
gives the coset tree for
$P=\langle\sigma\rangle\ast_3\langle\rho\rangle$.
The vertex stabilizers over the first kind of vertex are
exactly the conjugate copies of $\langle t\rangle$, indexed
by $P/\langle\sigma\rangle$; the other vertex and edge
stabilizers are trivial. The tree's universal property
identifies its kernel with their continuous free product.
The map in question sends each such factor isomorphically
to the corresponding actual inertia factor of Wingberg's
isomorphism. The free-product universal property gives its
inverse on the kernel, and the map on the quotient $P$ is
the identity. This proves (E55.5), including injectivity.

We made the choices of $D$ and $\rho$ in the full arithmetic
group. Its same elements now define a map in the reverse
direction from (E55.5). Projection back fixes all three
generators, so this map is a continuous section. Equation
(E55.6) makes it $c$-equivariant. Applying the same opened
theorem to $S$ and the set of all primes identifies the kernel
of this projection with the continuous free inertia product
over primes outside $S$. The $K(\pi,1)$ property for $S$
is also a conclusion of that theorem for $\{3\}\subset S$.
These statements specify the splitting and its exact
arithmetic domain; they do not discard the correction words
of E.39 for a different splitting.

\emph{Solving every unramified cocycle equation.}
For the fixed first cocycle, the coefficient module is
$J=(3,T-1)\mathbb Z_3[[T-1]]$ with operators
\[
 t(f)=T^{A_0}f,\quad
 \sigma(f)=qf(T^q),\quad
 \rho(f)=T^{B_0}f,\quad C(f)=-f(T^{-1}).                   \tag{E55.9}
\]
Normalize $v(c)=0$ by the order-two averaging argument.
Writing its other values as $(u,v_0,w)$, the complete
relations of (E55.5)--(E55.6) give exactly
\[
 \begin{split}
 u(T^{-1})&=T^{-A_0}u(T),\\
 v_0(T^{-1})&=-v_0(T),\\
 w(T^{-1})&=T^{-B_0}w(T),\\
 q u(T^q)&-(1+T^{A_0}+\cdots+T^{(q-1)A_0})u(T)
                                      =(T^{qA_0}-1)v_0(T).
 \end{split}                                                \tag{E55.10}
\]
The remaining geometric conjugations are precisely
$m\in J$ with $C(m)=m$, that is, anti-invariant series.
They send
\[
 (u,v_0,w)\longmapsto
 (u-(T^{A_0}-1)m,\ v_0-(q\psi_q-1)m,
                         \ w-(T^{B_0}-1)m),
 \quad\psi_q(f)=f(T^q).                                    \tag{E55.11}
\]
The first geometric coordinate must be zero here: preserving
the normalized first value at $c$ would otherwise change it
by twice that coordinate. Thus no conjugations have been lost.

Set $M=3^s$. The substitution $T\mapsto T^\epsilon$ is an
integral automorphism, and
\[
 (T^{A_0}-1)=(T^M-1),\qquad
 (T^{qA_0}-1)=(T^{A_0}-1)                                  \tag{E55.12}
\]
as ideals in the completed ring. The first equality follows
by the unit-exponent substitution, or by factoring its ratio
as a power series with constant term $\epsilon$ in $T^M-1$;
the ratio in the second equality has constant term $q$.
The quotient is $\mathbb Z_3[C_M]$. Reducing (E55.10) in it
shows that $\bar u$ is invariant under both inversion and
position multiplication by $q$. Since $v_3(q-1)=1$, the
closed subgroup generated by $q$ is $1+3\mathbb Z_3$;
together with $-1$ it acts transitively on positions of
each fixed three-power order. There are precisely $s+1$
orbits: zero and the orders $3,\ldots,3^s$.

The augmentation condition on $\bar u$ is exactly that its
total mass is in $3\mathbb Z_3$. This lattice has the basis
given by the images of $3$ and
$Q_d-\varphi(3^d)$ for $1\le d\le s$.
The factor $T^{A_0/2}$ is one in this quotient, so it
supplies a symmetric lift about $A_0/2$ without changing
any of these finite coordinates. This gives precisely
the first line of (E55.7).

For any such lift the numerator in its $\sigma$-value
vanishes modulo $T^{A_0}-1$, by these finite invariants.
Both ideal equalities in (E55.12) prove integral divisibility
by the displayed denominator. Inversion of its numerator
multiplies it by $T^{-qA_0}$; inversion of its denominator
multiplies it by $-T^{-qA_0}$. The quotient is therefore
anti-invariant, has augmentation zero, and lies in $J$.
This also proves all of (E55.10) for the proposed values.

Conversely take any solution of (E55.10). Its finite
coordinates determine a unique proposed $u_0$ with the
same image in $\mathbb Z_3[C_M]$. There is a unique
integral series
$m=(u-u_0)/(T^{A_0}-1)$. Inverting the numerator and
denominator shows $m(T^{-1})=-m(T)$, so $m\in J$ is an
allowed conjugation. It sends $u$ to $u_0$. The last
equation in (E55.10), in a domain, then fixes $v_0$
uniquely. No further nonzero conjugation fixes $u_0$.
Finally $T^{-B_0/2}w$ is invariant under inversion.
The change of coordinate $W=(T-T^{-1})/2$ identifies the
completed ring with $\mathbb Z_3[[W]]$, with inversion
$W\mapsto-W$. Its even subring is also
$\mathbb Z_3[[V]]$, since
$V=2(\sqrt{1+W^2}-1)$ has unit linear coefficient in $W^2$.
The augmentation restriction on $w$ gives exactly
$3\beta_3+Vf(V)$. This proves existence and uniqueness
of all coordinates in (E55.7).
The division maps used here are continuous on closed
principal ideals, by compactness and the domain property.
All remaining maps are continuous substitutions or finite
coefficient extractions. Thus this normal form is a
homeomorphism, including its stated integral topology.

\emph{All other primes and the exact global splitting.}
Let $L$ be the arithmetic kernel of the split projection
onto (E55.5). It acts trivially on $J_A$. Restriction to
$L$ of a second cocycle is therefore an additive continuous
homomorphism $\lambda:L\to J_A$. The splitting gives the
explicit reconstruction
\[
 v(l\,i(g))=\lambda(l)+v_S(g),\qquad
 \lambda(i(g)li(g)^{-1})=g\lambda(l),                         \tag{E55.13}
\]
where $i$ is the section just proved. This is a bijection
between actual cocycles, with $c$-equivariance imposed.
On the abelianized free inertia product, a different
prolongation of the same prime of $K_S(3)$ only conjugates
a factor by an element of $L$;
its image under $\lambda$ is unchanged. Consequently its
equivariant homomorphisms are exactly the compatible local
inertia values, with no unevaluated nonabelian correction.

At a rational prime $p\notin\{\ell,3\}$ the Tate-untwisted
action on positions is the affine map
$x\mapsto A(F_p)+px$. E.49 and E.51 prove directly that
its invariant integral measures in $J_A(-1)$ are zero
without an integral fixed point, and are precisely
$3\mathbb Z_3\delta_{x_p}$ with one. The fixed-point
condition is divisibility of the local Kummer class of
$\ell^{3a}3^{3b}$ by every power of three. Principal units
at $p$ are uniquely three-divisible; in its residue group
the condition says that the three-primary component is
one. Computing that component gives exactly (E55.2).
The representatives there are sufficient because changing
an exponent by $3^{h_p-1}$ changes its triple by a multiple
of $3^{h_p}$. In particular every $p$ with $h_p\le1$ is
allowed.

Normalize tame inertia at an allowed $p$ by
$\kappa_3(p)(t_p)=1$. Assign it
$\lambda_p(t_p)=3\delta_{x_p}$ and assign every other
base prime zero. Extend to all prolongations by the
equivariance in (E55.13), and then by the free-product
universal property. The fixed-point condition proves
independence of the chosen coset representative. At a
split rational prime the conjugate factor is prescribed
by $c$; at an inert prime the local rational Frobenius
condition already gives that descent. Thus $\lambda_p$
is an actual rational equivariant homomorphism.
Continuity of a family of such maps is exactly
$\beta_p\to0$ three-adically: augmentation forces this
condition at every finite coefficient exponent, and the
same condition makes the family uniformly zero at all
but finitely many primes in every finite quotient of $J$.
This is the continuous-bundle topology of E.39.

There is a small toric normalization to record. The
cocycle reconstructed with $v_S=0$ from $\lambda_p$
has scaled second coordinate
\[
             \kappa_3(p)-\kappa_3(p)(\rho)\kappa_3(3).
 \tag{E55.14}
\]
Indeed it has valuation one at $p$, zero at every other
prime except possibly three, and vanishes on $\rho$.
At $\ell$ its whole decomposition restriction is zero.
Also $\kappa_3(p)(\sigma)=0$: the chosen $c$ fixes
$\sigma$, whereas its Kummer action is negation and its
value at $c$ is zero. The rational completed Kummer
coordinates of E.16 now prove (E55.14), including the
coefficient at three. Add the unramified cocycle (E55.7)
with $\beta_3=\kappa_3(p)(\rho)$ and every other coordinate
zero. Denote the result by $v_p$. It has scaled second
coordinate exactly $\kappa_3(p)$.

Every full global cocycle has unique residues
$3\beta_p\delta_{x_p}$. Subtract the convergent sum
$\sum_p\beta_p v_p$; (E55.13) shows that the remainder
is unramified outside $\{\ell,3\}$, and (E55.7) gives its
unique remaining coordinates. The terms used to normalize
$v_p$ have coefficients in $\mathbb Z_3$, so the same
restrictedness proves convergence of that normalization
as well. All constructions and their inverses are
continuous at every finite coefficient and positional
level. This proves (E55.3) and (E55.4) as actual topological
group coordinates. It also proves three-torsion-freeness.

Finally a global section of $H$ with the specified first
class is exactly this twisted additive cocycle modulo
the conjugations used above, by the multiplication law
of E.46. Its finite images over $K$ are three-groups,
so passing to the relative pro-three arithmetic quotient
preserves all these sections and all their conjugacies.
The maps from the intrinsic Fermat groups to $H$ retain
their previously computed kernels. The classification
here applies to this specified target and first family.

For $\ell=2$ and $b=0$, the comparison with the norm
coordinates of E.52 is explicit through their common
section cocycle. Apply the finite Kummer--Shapiro maps
of E.46 to a norm system, evaluate its cocycle on
$t,\sigma,\rho,c$, and normalize its value at $c$.
Its image in $\mathbb Z_3[C_{3^s}]$ supplies the
coefficients of $3$ and $Q_d-\varphi(3^d)$; division
by $T^{A_0}-1$ then gives the unique conjugator in
(E55.11), and the remaining $\rho$-value gives $f$.
For the inverse, reconstruct that actual cocycle by
(E55.7) and (E55.13), apply the inverse Shapiro maps,
and extract the valuations and successive unit
coordinates by (E52.10)--(E52.13). Both composites
are the identity because each of those component
maps was proved bijective. This specifies the
arithmetic homeomorphism between the two displayed
parameter groups, including their finite coordinate
blocks and their norm transitions.
\end{proof}
